\noindent\fbox{\begin{minipage}{\textwidth}
\paragraph{Theorem 2:} Let
\[
N_k=\#\{(x,y)\in \mathbb Z^2:\ x^3+y^3=x^2+kxy+y^2\}.
\]
If \(k\neq -1,2\), then
\[
N_k\le 4\,\tau\!\left(\left|(k-2)^2(k+1)\right|\right),
\]
where \(\tau(m)\) denotes the number of positive divisors of \(m\).

Moreover, for \(k=-1\) and \(k=2\), the equation has infinitely many integer solutions.
\end{minipage}}
\\
\\
\textbf{Proof.}
Consider the equation
\[
x^3+y^3=x^2+kxy+y^2.
\]
Set
\[
s=x+y,\qquad p=xy.
\]
Then
\[
x^3+y^3=(x+y)^3-3xy(x+y)=s^3-3ps,
\]
and
\[
x^2+kxy+y^2=(x+y)^2+(k-2)xy=s^2+(k-2)p.
\]
Hence
\[
s^3-3ps=s^2+(k-2)p,
\]
so
\[
s^3-s^2=p(3s+k-2),
\]
that is,
\[
s^2(s-1)=p(3s+k-2).
\]
Let
\[
D=3s+k-2.
\]
Then
\[
D\mid s^2(s-1).
\]

We first note that if \(D=0\), then
\[
s^2(s-1)=0,
\]
so \(s=0\) or \(s=1\). Since \(D=3s+k-2\), this gives \(k=2\) or \(k=-1\). Therefore, for
\[
k\neq -1,2,
\]
we must have
\[
D\neq 0.
\]

We now prove that
\[
|D|\mid |(k-2)^2(k+1)|.
\]

Let \(q^a\parallel D\), where \(q\) is a prime. Since
\[
D\mid s^2(s-1)
\]
and
\[
\gcd(s,s-1)=1,
\]
the prime \(q\) divides exactly one of \(s\) and \(s-1\).

\medskip

\textit{Case 1: \(q\mid s-1\).}

Since \(q\nmid s\), from \(q^a\mid s^2(s-1)\) we get
\[
q^a\mid (s-1).
\]
But also \(q^a\mid D\), hence
\[
q^a\mid D-3(s-1)=k+1.
\]

\medskip

\textit{Case 2: \(q\mid s\).}

Write
\[
q^b\parallel s.
\]
Since \(q\nmid s-1\), from \(q^a\mid s^2(s-1)\) we obtain
\[
a\le 2b.
\]
Also, because \(q^a\mid D\) and \(q^b\mid s\), we have
\[
q^{\min(a,b)}\mid \gcd(D,s).
\]
But
\[
D-3s=k-2,
\]
so
\[
q^{\min(a,b)}\mid (k-2).
\]
If \(a\le b\), then \(q^a\mid (k-2)\). If \(a>b\), then \(q^b\mid (k-2)\), and since \(a\le 2b\), it follows that
\[
q^a\mid (k-2)^2.
\]
Thus in all cases,
\[
q^a\mid (k-2)^2.
\]

Therefore, every prime power divisor of \(D\) divides either \(k+1\) or \((k-2)^2\), and so
\[
|D|\mid |(k-2)^2(k+1)|.
\]

Now let
\[
M=(k-2)^2(k+1).
\]
Every integer solution \((x,y)\) gives rise to a divisor \(D\) of \(M\), and \(s=x+y\) is then uniquely determined by
\[
s=\frac{D-k+2}{3}.
\]
After that, \(p=xy\) is uniquely determined from
\[
p=\frac{s^2(s-1)}{D}.
\]

Finally, \(x\) and \(y\) are the roots of
\[
t^2-st+p=0.
\]
For fixed \(s\) and \(p\), this quadratic gives at most two ordered integer pairs \((x,y)\), namely \((x,y)\) and \((y,x)\).

Since \(M\neq 0\) for \(k\neq -1,2\), the number of integer divisors of \(M\) is
\[
2\tau(|M|).
\]
Hence
\[
N_k\le 2\cdot 2\tau(|M|)=4\tau(|M|).
\]
Therefore
\[
N_k\le 4\,\tau\!\left(\left|(k-2)^2(k+1)\right|\right).
\]

\bigskip

It remains to note the exceptional cases from introduction.
This completes the proof.